\section*{Capítulo \ref{ch:b3:quantum-model-systems} --- Mecánica cuántica para químicos: sistemas modelo}

\begin{solution}{exo:b3:quantum-model-systems:1}
$E_1 = h^2/8m_eL^2$. (a) $L = \qty{1.00}{nm}$: $E_1 = \qty{6.02e-20}{J} =
\qty{0.376}{eV}$; $\Delta E_{12} = 3E_1$ y $\lambda = hc/3E_1 =
\qty{1099}{nm}$, en el infrarrojo cercano. (b) $L = \qty{0.50}{nm}$: $E_1$ es cuatro
veces mayor, \qty{2.41e-19}{J} $= \qty{1.50}{eV}$, y $\lambda =
\qty{275}{nm}$, en el ultravioleta.
\end{solution}

\begin{solution}{exo:b3:quantum-model-systems:2}
$n_x^2 + n_y^2 + n_z^2 = 3$ (1,1,1): $g = 1$; $6$ (2,1,1 y permutaciones):
$g = 3$; $9$ (2,2,1): $g = 3$; $11$ (3,1,1): $g = 3$; $12$ (2,2,2): $g = 1$.
\end{solution}

\begin{solution}{exo:b3:quantum-model-systems:3}
$[\hat x^2,\hat p]f = -\iu\hbar x^2f' + \iu\hbar(x^2f)' = 2\iu\hbar xf$, luego
$[\hat x^2,\hat p] = 2\iu\hbar\hat x$. $[\hat p^2,\hat x]f = -\hbar^2[(xf)'' -
xf''] = -2\hbar^2f'$, luego $[\hat p^2,\hat x] = -2\iu\hbar\hat p$.
\end{solution}

\begin{solution}{exo:b3:quantum-model-systems:4}
$E_J/hc = BJ(J+1)$: 0, 21.19, 63.56 y \qty{127.12}{cm^{-1}}, con $g = 1$,
3, 5, 7. $I = h/(8\pi^2cB) = \qty{2.643e-47}{kg.m^2}$ (con $c$ en
\unit{cm/s}).
\end{solution}

\begin{solution}{exo:b3:quantum-model-systems:5}
$|\psi_n|^2$ es simétrica respecto de $L/2$, luego $\langle x\rangle = L/2$. Con
$\sin^2u = \frac12(1 - \cos2u)$ y dos integraciones por partes,
$\langle x^2\rangle = L^2\big(\frac13 - \frac{1}{2n^2\pi^2}\big)$. Para $n = 1$,
$\langle x^2\rangle = 0.283L^2$, menos que el valor clásico $L^2/3$: el estado
fundamental se acumula en el centro. Al crecer $n$ se recupera el valor clásico.
\end{solution}

\begin{solution}{exo:b3:quantum-model-systems:6}
$1\,\unit{cm^{-1}} = \qty{0.011963}{kJ/mol}$. EPC(\ce{H2}) $= 2200.6 \times
0.011963 = \qty{26.33}{kJ/mol}$, EPC(\ce{D2}) $= \qty{18.63}{kJ/mol}$; diferencia:
\qty{7.69}{kJ/mol}. Misma \omterm{def:b3:quantum-model-systems:oscillator}{constante de fuerza}, luego $\tilde\omega \propto \mu^{-1/2}$
y la razón esperada es $\sqrt{\mu_D/\mu_H} = 1.4137$ (masas atómicas
2.01410 y 1.00783); la razón medida es 1.4127, a un 0.07\,\% (la
pequeña diferencia procede de los electrones, que no siguen a los núcleos
perfectamente).
\end{solution}

\begin{solution}{exo:b3:quantum-model-systems:7}
$\int_0^Lx^2(L - x)^2\,\dd x = L^5/30$, luego $N = \sqrt{30/L^5}$. Entonces
\[
\langle\psi_1|\phi\rangle = \sqrt{\tfrac2L}\sqrt{\tfrac{30}{L^5}}\int_0^Lx(L -
x)\sin\tfrac{\pi x}{L}\,\dd x = \frac{\sqrt{60}}{L^3}\cdot\frac{4L^3}{\pi^3} =
\frac{4\sqrt{60}}{\pi^3} = 0.99928 .
\]
La parábola es en un 99.86\,\% estado fundamental (el cuadrado del solapamiento).
\end{solution}

\begin{solution}{exo:b3:quantum-model-systems:8}
$\langle f|\hat pg\rangle = -\iu\hbar\int f^*g'\,\dd x = -\iu\hbar[f^*g] +
\iu\hbar\int f^{*\prime}g\,\dd x = \int(-\iu\hbar f')^*g\,\dd x = \langle\hat
pf|g\rangle$, pues el corchete se anula en los extremos. Para $\hat x\hat p$, usando que $\hat x$ y $\hat p$ son ambos \omterm{def:b3:quantum-model-systems:hermitian}{hermíticos}:
$\langle f|\hat x\hat pg\rangle = \langle\hat xf|\hat pg\rangle = \langle\hat p\hat
xf|g\rangle$, y $\hat p\hat x = \hat x\hat p - \iu\hbar \ne \hat x\hat p$: no es
\omterm{def:b3:quantum-model-systems:hermitian}{hermítico} (su forma simetrizada $\frac12(\hat x\hat p + \hat p\hat x)$ sí lo es).
\end{solution}

\begin{solution}{exo:b3:quantum-model-systems:9}
$L = 6 \times 139.7 = \qty{838}{pm}$, $N = 6$: $\lambda = 8m_ecL^2/(7h) =
\qty{331}{nm}$. La absorción está en el ultravioleta: el hexatrieno es
incoloro.
\end{solution}

\begin{solution}{exo:b3:quantum-model-systems:10}
$\kappa = \sqrt{2m(V_0 - E)}/\hbar$ con $V_0 - E = \qty{0.20}{eV}$:
$\kappa_H = \qty{9.82e10}{m^{-1}}$, $\kappa_D = \qty{1.388e11}{m^{-1}}$. Los
prefactores se cancelan en la razón: $T_H/T_D = \eu^{2a(\kappa_D - \kappa_H)} =
\eu^{4.06} = 58$, el valor exacto: la barrera es gruesa para ambos
($\kappa a = 4.9$ y 6.9).
\end{solution}

\begin{solution}{exo:b3:quantum-model-systems:11}
Niveles $E_m = m^2\hbar^2/2m_eR^2$: $m = 0$ aloja dos electrones, $m = \pm1$
cuatro. El nivel lleno más alto es $|m| = 1$, el vacío más bajo $|m| = 2$: $\Delta E =
3\hbar^2/2m_eR^2 = \qty{9.38e-19}{J}$, $\lambda = \qty{212}{nm}$, en el
ultravioleta, como la fuerte absorción del benceno.
\end{solution}

\begin{solution}{exo:b3:quantum-model-systems:12}
$\langle r\rangle = \frac{4\pi}{\pi a^3}\int_0^\infty r^3\eu^{-2r/a}\dd r =
\frac{4}{a^3}\cdot\frac{3!}{(2/a)^4} = \frac32a$, \qty{79.4}{pm} para $a = a_0$.
$\langle 1/r\rangle = \frac{4}{a^3}\cdot\frac{1!}{(2/a)^2} = \frac1a$, luego
$\langle V\rangle = -e^2/(4\pi\varepsilon_0a)$. Como $E_1 =
-e^2/(8\pi\varepsilon_0a)$ (a partir de $a = 4\pi\varepsilon_0\hbar^2/\mu e^2$),
$\langle V\rangle = 2E_1$, y la energía cinética media es $-E_1$.
\end{solution}

\begin{solution}{pb:b3:quantum-model-systems:1}
\textbf{1.} Cada uno de los $j$ átomos aporta un orbital p; los $j - 1$ carbonos y
un nitrógeno aportan un electrón cada uno y el otro nitrógeno su par no enlazante,
ya que la carga del catión está repartida por la cadena: $N = j + 1$, es decir, 6,
8 y 10.
\textbf{2.} Seis electrones llenan $n = 1$, 2, 3: HOMO $n = 3$, LUMO $n = 4$.
\textbf{3.} $\Delta E = E_{N/2+1} - E_{N/2} = (N+1)h^2/8mL^2$ y $\lambda =
hc/\Delta E = 8mcL^2/h(N+1)$.
\textbf{4.} $L_0 = 4l$, $6l$, $8l$: 559, 838 y \qty{1118}{pm}.
\textbf{5.} 147, 257 y \qty{374}{nm}.
\textbf{6.} Todas demasiado cortas (por factores de 3.6, 2.3 y 1.9): la caja del modelo es demasiado
pequeña, ya que $\lambda \propto L^2$; los electrones se mueven sobre una distancia
mayor que la cadena comprendida entre los nitrógenos.
\textbf{7.} Con el origen en el centro, $\psi_n \propto \cos(n\pi u/L)$ para
$n$ impar y $\sin(n\pi u/L)$ para $n$ par ($u = x - L/2$): funciones par e impar
de $u$.
\textbf{8.} Si $n + n'$ es par, $\psi_n$ y $\psi_{n'}$ tienen la misma paridad;
su producto multiplicado por $u$ es impar y su integral sobre un intervalo simétrico
es nula.
\textbf{9.} El HOMO $N/2$ y el LUMO $N/2 + 1$ tienen $n + n' = N + 1$, impar: siempre
permitida.
\textbf{10.} Segundo colorante: HOMO 4, LUMO 5. $4 \to 6$: suma par, prohibida.
$3 \to 5$: suma par, prohibida.
\textbf{11.} La siguiente permitida desde $n = 4$ es $4 \to 7$, con un
$\Delta E$ mayor en un factor $(49 - 16)/(25 - 16) = 3.67$: a $\lambda/3.67$,
\qty{164}{nm} para el segundo colorante, en pleno ultravioleta.
\textbf{12.} Solo depende de la simetría del potencial respecto del centro de la
cadena, que comparten los colorantes simétricos reales.
\textbf{13.} $E = hc/\lambda = \qty{3.294e-19}{J} = \qty{2.06}{eV}$.
\textbf{14.} $hc\tilde\omega_e = \qty{0.371}{eV}$.
\textbf{15.} $2hcB = \qty{4.21e-22}{J} = \qty{2.63}{meV}$.
\textbf{16.} $kT = \qty{25.7}{meV}$.
\textbf{17.} Solo los niveles de rotación ($2hcB \ll kT$); las excitaciones vibracionales ($15kT$) y
electrónicas ($80kT$) prácticamente no están pobladas.
\textbf{18.} Electrónicas: visible y ultravioleta; vibracionales: infrarrojo;
rotacionales: microondas (e infrarrojo lejano).
\textbf{19.} $L = \sqrt{\lambda h(N+1)/8mc}$ con $N = 8$: $L = \qty{1283}{pm}$.
\textbf{20.} $\delta = (1283 - 838)/2 = \qty{222}{pm}$.
\textbf{21.} $L = 559 + 445 = \qty{1003}{pm}$, $\lambda = \qty{474}{nm}$, un 9.5\,\%
por debajo de 524 nm.
\textbf{22.} $L = 1118 + 445 = \qty{1562}{pm}$, $\lambda = \qty{732}{nm}$, un 2.9\,\%
por encima de 711 nm.
\textbf{23.} Hacia los anillos aromáticos que llevan los nitrógenos: el sistema
$\pi$ no termina en los átomos de nitrógeno, y la caja efectiva incluye
parte de cada anillo.
\textbf{24.} $\delta/l = 222/139.7 = 1.6$: \textbf{la caja se prolonga unos
\qty{222}{pm}, enlace y medio, más allá de cada nitrógeno}, y con esa sola
longitud el modelo reproduce los tres colores con un error inferior al 10\,\%.
\end{solution}
